\section{Discussion and Limits}

This note is deliberately careful about what it does not prove.

The cleanest way to state the gap is as a three-part checklist. First, the theorem here assumes additivity on the full effect space rather than only on orthogonal projections. Second, it stays in the real symmetric setting rather than the full complex Hilbert-space statement in \citet{gleason1957}. Third, it is a finite-dimensional theorem. In that sense the note sits much closer to the effect-space perspective emphasized by \citet{busch2003} and discussed cleanly by \citet{wright2019} than to the original projection-only theorem.

It also does not claim that the effect-space theorem proved here is new as mathematics. The mathematical contribution is expository and structural: the point is to make the route to the trace form feel inevitable to a newcomer rather than miraculous on first contact.

What the note does contribute is a more teachable route. The proof starts from a one-variable bounded-additivity lemma, turns that into a sphere warmup that can be explained with squared heights, and only then moves to real matrices. That progression matters if the goal is pedagogy, onboarding, or conceptual demystification rather than maximal theorem strength.

There is also a limit to the computational layer. The generated checks confirm that the reconstruction formulas and state-rule identities behave correctly in random sampled examples. They do not strengthen the theorem mathematically. In a theoretical paper, the proof remains primary.

The clearest next step would be a follow-on note that closes one part of the gap at a time: first by recovering more of the projection-only setting, then by clarifying which parts of the argument survive unchanged in the complex case, while preserving as much of the present accessibility as possible.
